Все графы конечные и простые (без петель и кратных рёбер), если не оговорено иное;
\(G=(V,E)\), \(n=|V|\), \(m=|E|\).

\subsection*{Основные обозначения}
{\renewcommand{\arraystretch}{1.22}
\begin{tabular}{@{}p{0.30\linewidth}@{\hspace{1em}}p{\dimexpr0.70\linewidth-1em\relax}@{}}
\textbf{Обозначение} & \textbf{Значение} \\
\hline
\(N(v)\), \(N(A)\) & Соседи вершины; \(N(A)=\bigcup_{a\in A}N(a)\) \\
\(G-S\), \(G[S]\) & Удаление множества вершин; индуцированный подграф \\
\(H\le G\) & \(H\) — индуцированный подграф \(G\) \\
\(\overline G\) & Дополнение: рёбра ровно между несмежными в \(G\) вершинами \\
\(V(M)\) & Множество концов рёбер \(M\subseteq E\) \\
\(M\mathbin{\triangle}M'\) & Симметрическая разность (в лекциях \(M\oplus M'\)) \\
\(\odd(G)\), \(\comp(G)\) & Число нечётных компонент; число всех компонент связности \\
\(K_n\), \(K_{p,q}\), \(C_n\), \(P_n\) & Полный, полный двудольный граф, цикл и путь на \(n\) вершинах
\end{tabular}\par}

\subsection*{Множества вершин}
\begin{definition}[Независимое множество]
\(I\subseteq V\) \emph{независимо}, если никакие две его вершины не смежны.
\(\alpha(G)=\max|I|\) — \emph{число независимости}.
\end{definition}
\begin{definition}[Вершинное покрытие]
\(C\subseteq V\) — \emph{вершинное покрытие}, если у каждого ребра хотя бы один конец в \(C\).
\(\beta(G)=\min|C|\) (в лекциях \(\mathrm{vc}(G)\)).
\end{definition}
\begin{definition}[Клика]
\(K\subseteq V\) — \emph{клика}, если \(G[K]\) полный. \(\omega(G)=\max|K|\) — \emph{кликовое число};
\(\chi(G)\) — хроматическое число. Заметим \(\alpha(G)=\omega(\overline G)\).
\end{definition}

\subsection*{Множества рёбер}
\begin{definition}[Паросочетание]
\(M\subseteq E\) — \emph{паросочетание}, если никакие два его ребра не имеют общего конца.
\(\alpha'(G)=\max|M|\) (в лекциях \(\mathrm{mm}(G)\)).
\(M\) \emph{покрывает} вершину \(v\), если \(v\in V(M)\); остальные вершины \emph{свободны}.
\end{definition}
\begin{definition}[Рёберное покрытие]
\(F\subseteq E\) — \emph{рёберное покрытие}, если каждая вершина — конец ребра из \(F\).
\(\beta'(G)=\min|F|\); существует лишь при \(\delta(G)>0\).
\end{definition}
\begin{definition}[Совершенное и почти совершенное паросочетание]
Паросочетание \emph{совершенно}, если покрывает все вершины (тогда \(|M|=n/2\)),
и \emph{почти совершенно}, если покрывает все вершины, кроме одной.
\end{definition}
\begin{definition}[Чередующийся и дополняющий путь]
Простой путь \emph{\(M\)-чередующийся}, если его рёбра поочерёдно лежат и не лежат в \(M\).
\emph{\(M\)-дополняющий} путь — \(M\)-чередующийся путь положительной длины с двумя
свободными концами; рёбер не из \(M\) в нём на одно больше.
\end{definition}
\begin{tcolorbox}[colback=cbOrangeBack,colframe=cbOrange!45!cbGray,
  borderline west={1.6pt}{0pt}{cbOrange}]
\textbf{Соглашение.} «Максимальное» паросочетание (независимое множество) — наибольшее
по размеру. Не расширяемое добавлением элемента называем \emph{максимальным по
включению}; оно может быть меньше (рис.~\ref{fig:notation-matchings}).
\end{tcolorbox}
\begin{graphillustration}{Путь $P_4$: среднее ребро — паросочетание, максимальное по включению (ни одно ребро не добавить), но не наибольшее. Два крайних ребра — совершенное паросочетание.}{fig:notation-matchings}
\begin{tikzpicture}[font=\small]
\begin{scope}
\node[gi free] (a) at (0,0) {$a$};
\node[gi vertex] (b) at (1.3,0) {$b$};
\node[gi vertex] (c) at (2.6,0) {$c$};
\node[gi free] (d) at (3.9,0) {$d$};
\draw[gi edge] (a) -- (b);
\draw[gi match] (b) -- (c);
\draw[gi edge] (c) -- (d);
\node[gi note] at (1.95,-0.8) {$M=\{bc\}$: максимально по включению,\\$a,b,c,d$ — дополняющий путь};
\end{scope}
\begin{scope}[shift={(6.5,0)}]
\node[gi vertex] (a) at (0,0) {$a$};
\node[gi vertex] (b) at (1.3,0) {$b$};
\node[gi vertex] (c) at (2.6,0) {$c$};
\node[gi vertex] (d) at (3.9,0) {$d$};
\draw[gi match] (a) -- (b);
\draw[gi edge] (b) -- (c);
\draw[gi match] (c) -- (d);
\node[gi note] at (1.95,-0.8) {$M'=\{ab,cd\}$: совершенное,\\$|M'|=\alpha'(P_4)=2$};
\end{scope}
\end{tikzpicture}
\end{graphillustration}


\subsection*{Раскраски рёбер}
\begin{definition}[Рёберная раскраска]
Раскраска рёбер \(K_t\) в \(k\) цветов — функция \(\varphi\colon E(K_t)\to\{1,\ldots,k\}\).
При \(k=2\) она равносильна графу \(G=(V,\varphi^{-1}(1))\): рёбра цвета 2 образуют
\(\overline G\), клики цвета 2 — независимые множества \(G\).
\end{definition}
\textbf{На рисунках:} рёбра паросочетания — толстые синие; свободные вершины — двойной
серый контур; цвета 1, 2, 3 — сплошная синяя, штриховая оранжевая, пунктирная зелёная линии.

\subsection*{Исходные факты}
\begin{lemma}[О рукопожатиях]
\(\sum_{v}d(v)=2m\); число вершин нечётной степени чётно.
\end{lemma}
\begin{lemma}[Рёбра и компоненты]
В любом графе \(m\ge n-\comp(G)\).
\end{lemma}
\begin{lemma}[Клика и раскраска]
\(\chi(G)\ge\omega(G)\).
\end{lemma}
